\documentclass[10pt,a4paper,final]{article} %draft

\usepackage[english, russian]{babel}

\usepackage{geometry}
%\usepackage[a4paper,left=2cm,right=2cm,top=2cm,bottom=2cm,bindingo ffset=0cm]{geometry}

\usepackage[T2A]{fontenc}
\usepackage[utf8]{inputenc}

\usepackage[utf8]{inputenc}

\usepackage[final]{pdfpages}

\usepackage{textcomp,enumitem}

\usepackage{amsmath,amsthm,amssymb}

\usepackage{longtable}
\usepackage{array}
\usepackage{adjustbox}
\usepackage{pdflscape}  % Лучше, чем lscape, для PDF
\usepackage{fancyhdr}



\usepackage{listings}
\usepackage{xcolor}
\usepackage{caption}

\definecolor{codegray}{rgb}{0.95,0.95,0.95}
\definecolor{codepurple}{rgb}{0.58,0,0.82}
\definecolor{codegreen}{rgb}{0,0.5,0} % зеленый цвет
\definecolor{codeblue}{rgb}{0,0,0.5} % синий цвет
\definecolor{codestring}{rgb}{0.64,0.08,0.08} % цвет строк


\lstset{
	language=C++,
	extendedchars=true,
	literate=
	{а}{{\selectfont\char224}}1
	{б}{{\selectfont\char225}}1
	{в}{{\selectfont\char226}}1
	{г}{{\selectfont\char227}}1
	{д}{{\selectfont\char228}}1
	{е}{{\selectfont\char229}}1
	{ж}{{\selectfont\char230}}1
	{з}{{\selectfont\char231}}1
	{и}{{\selectfont\char232}}1
	{й}{{\selectfont\char233}}1
	{к}{{\selectfont\char234}}1
	{л}{{\selectfont\char235}}1
	{м}{{\selectfont\char236}}1
	{н}{{\selectfont\char237}}1
	{о}{{\selectfont\char238}}1
	{п}{{\selectfont\char239}}1
	{р}{{\selectfont\char240}}1
	{с}{{\selectfont\char241}}1
	{т}{{\selectfont\char242}}1
	{у}{{\selectfont\char243}}1
	{ф}{{\selectfont\char244}}1
	{х}{{\selectfont\char245}}1
	{ц}{{\selectfont\char246}}1
	{ч}{{\selectfont\char247}}1
	{ш}{{\selectfont\char248}}1
	{щ}{{\selectfont\char249}}1
	{ъ}{{\selectfont\char250}}1
	{ы}{{\selectfont\char251}}1
	{ь}{{\selectfont\char252}}1
	{э}{{\selectfont\char253}}1
	{ю}{{\selectfont\char254}}1
	{я}{{\selectfont\char255}}1
	{А}{{\selectfont\char192}}1
	{Б}{{\selectfont\char193}}1
	{В}{{\selectfont\char194}}1
	{Г}{{\selectfont\char195}}1
	{Д}{{\selectfont\char196}}1
	{Е}{{\selectfont\char197}}1
	{Ж}{{\selectfont\char198}}1
	{З}{{\selectfont\char199}}1
	{И}{{\selectfont\char200}}1
	{Й}{{\selectfont\char201}}1
	{К}{{\selectfont\char202}}1
	{Л}{{\selectfont\char203}}1
	{М}{{\selectfont\char204}}1
	{Н}{{\selectfont\char205}}1
	{О}{{\selectfont\char206}}1
	{П}{{\selectfont\char207}}1
	{Р}{{\selectfont\char208}}1
	{С}{{\selectfont\char209}}1
	{Т}{{\selectfont\char210}}1
	{У}{{\selectfont\char211}}1
	{Ф}{{\selectfont\char212}}1
	{Х}{{\selectfont\char213}}1
	{Ц}{{\selectfont\char214}}1
	{Ч}{{\selectfont\char215}}1
	{Ш}{{\selectfont\char216}}1
	{Щ}{{\selectfont\char217}}1
	{Ъ}{{\selectfont\char218}}1
	{Ы}{{\selectfont\char219}}1
	{Ь}{{\selectfont\char220}}1
	{Э}{{\selectfont\char221}}1
	{Ю}{{\selectfont\char222}}1
	{Я}{{\selectfont\char223}}1,
	%	{|}{{\textbar}}1,
	%	{||}{{\textbar\textbar}}1
	%	{\&}{{\string&}}1
	%	{==}{{\string==}}1
	%	{\string\\}{{\string\\}}1,
	numbers=left,
	stepnumber=1,
	firstnumber=1,
	numberstyle=\tiny,
	basicstyle=\ttfamily\footnotesize,
	breakatwhitespace=false,
	breaklines=true,
	captionpos=b,
	keepspaces=true,
	showspaces=false,
	showstringspaces=false,
	showtabs=false,
	tabsize=2,
	frame=single,
	keywordstyle=\color{codeblue}, % ключевые слова
	commentstyle=\color{codegreen}, % комментарии
	stringstyle=\color{codestring}, % строки
	%identifierstyle=\color{green}, % стиль для переменных
	backgroundcolor=\color{codegray},
}
%\usepackage{fancyvrb}

%\usepackage{fancyhdr}

%\usepackage{upgreek}

%\usepackage{tipa}

%\usepackage{tikz}

%\usepackage{graphicx}

%\usepackage{pgfplots}

\usepackage{indentfirst}

%\usepackage{gensymb}

%\usepackage{amssymb}

%\usepackage{pdfpages}

\usepackage[unicode, pdftex, colorlinks, urlcolor=blue]{hyperref}

%\usepackage[T2A]{fontenc}


\usepackage{tabularx}
\usepackage{multirow}

\usepackage{booktabs} % для стильных линий таблицы



\usepackage{graphics}

\linespread{1.5}

\pagestyle{plain}

%\usepackage{listings}
%\usepackage{moreverb}

%\setlist[enumerate,itemize]{leftmargin=1.2cm} %отступ в перечислениях

\hypersetup{
	colorlinks=true,
	linkcolor=black,
	%allcolors=[RGB]{0 0 0}}
	}
	%\lstloadlanguages{ [LaTeX] TeX}

	\lstloadlanguages{ [LaTeX] TeX}



	\textheight=24cm
	\textwidth=16cm
	\oddsidemargin=0pt
	\topmargin=-1.5cm
	\parindent=24pt
	\parskip=0pt
	\tolerance=2000
	\flushbottom

	%\usepackage[font=scriptsize]{caption}
	\usepackage[labelsep=period]{caption}

	\begin{document}
\thispagestyle{empty}

\begin{center}
	{\Large МИНОБРНАУКИ РОССИИ}\\
	~\\
	{\large ФЕДЕРАЛЬНОЕ ГОСУДАРСТВЕННОЕ БЮДЖЕТНОЕ ОБРАЗОВАТЕЛЬНОЕ УЧРЕЖДЕНИЕ ВЫСШЕГО ПРОФЕССИОНАЛЬНОГО ОБРАЗОВАНИЯ}\\
	~\\
	{\Large \bf <<САНКТ-ПЕТЕРБУРГСКИЙ ПОЛИТЕХНИЧЕСКИЙ УНИВЕРСИТЕТ ПЕТРА ВЕЛИКОГО>>}\\
	~\\
	{\large Институт компьютерных наук и кибербезопасности}\\
	{\large Высшая школа технологий искусственного интеллекта}\\
	{\large Направление 02.03.01 Математика и компьютерные науки}\\
	~\\
	~\\
	~\\
	~\\
	{\Large \bf Lessons and Exams Timetable Bot}\\
		{\Large \bf Архитектура}\\
	{\Large Гладков И.А., Саржан А.А., Смирнов С.Р., Фролов К.А.}\\
	~\\
	~\\


	~\\
	~\\
	~\\
	~\\
	~\\
	~\\
	~\\

	~\\
	~\\
	~\\

	~\\
	~\\
	~\\

	~\\
	~\\
	~\\




	{\large Санкт-Петербург, 2026}
\end{center}
\newpage

\tableofcontents

\newpage

\begin{landscape}
\begin{figure}[H]
\thispagestyle{empty}
    \centering
    \includegraphics[scale=0.5]{arc.png}
    \caption{Архитектура}
    \vspace{30mm}4
    \label{arc}
\end{figure}
\end{landscape}
\newpage

\section{Компонентная диаграмма}

На компонентной диаграмме представлена архитектура Telegram-бота
Lessons and Exams Timetable Bot. Система построена как серверное
приложение на Java и Spring WebFlux. Внутри сервера выделены следующие
логические слои: входной API layer, бизнес-модули, репозитории и адаптеры
для взаимодействия с внешними системами.

Основной пользовательский сценарий начинается с отправки команды или
файла через Telegram. Запрос поступает в слой API, затем передаётся в
\texttt{MessageHandlerModule}, который определяет тип команды и направляет
её в соответствующий бизнес-модуль. Бизнес-модули работают с данными
через репозитории и внешние API-адаптеры. Данные пользователей,
расписаний, задач, напоминаний и учебных сессий сохраняются в MongoDB.
Для асинхронной обработки напоминаний используется RabbitMQ. Диаграмма представлена на рисунке \ref{arc}.


\subsection{Сервер}

\textbf{Роль:} центральный компонент системы.

\textbf{Функционал:}
\begin{itemize}
    \item принимает пользовательские команды, сообщения и файлы через Telegram;
    \item принимает HTTP-запросы администратора;
    \item обрабатывает служебный endpoint \texttt{/healthcheck};
    \item координирует взаимодействие между API layer, бизнес-модулями,
    репозиториями и внешними API;
    \item обеспечивает работу с MongoDB, RabbitMQ, LLM API, календарными API
    и внешним API расписания учебного заведения.
\end{itemize}

\subsection{Пользователь (User)}

\textbf{Роль:} основной конечный пользователь Telegram-бота.

\textbf{Функционал:}
\begin{itemize}
    \item взаимодействует с системой через Telegram;
    \item отправляет команды для просмотра расписания, управления задачами,
    настройки напоминаний и получения LLM-рекомендаций;
    \item может загружать файл расписания в формате CSV или iCal;
    \item идентифицируется автоматически по Telegram \texttt{chatId} или
    \texttt{telegramId};
    \item может выполнить проверку состояния системы через endpoint
    \texttt{/healthcheck}, если такой endpoint доступен публично.
\end{itemize}

\subsection{Администратор (Admin)}

\textbf{Роль:} пользователь с расширенными правами.

\textbf{Функционал:}
\begin{itemize}
    \item взаимодействует с системой через HTTP API;
    \item получает API-ключ или Bearer-токен через endpoint \texttt{/auth};
    \item обращается к административному endpoint \texttt{/users};
    \item может получать список пользователей системы;
    \item может выполнять проверку состояния сервера через endpoint
    \texttt{/healthcheck}.
\end{itemize}

\subsection{API layer}

\textbf{Роль:} входной слой приложения.

\textbf{Функционал:}
\begin{itemize}
    \item принимает входящие запросы от пользователей и администратора;
    \item преобразует внешние запросы в формат, удобный для бизнес-модулей;
    \item передаёт команды, данные пользователя, файлы и HTTP-запросы в
    соответствующие модули.
\end{itemize}

\subsubsection{API (Telegram message handlers)}

\textbf{Роль:} обработка входящих Telegram-сообщений и команд.

\textbf{Функционал:}
\begin{itemize}
    \item принимает команды пользователя, например \texttt{/start},
    \texttt{/today}, \texttt{/week}, \texttt{/add\_task}, \texttt{/advisor};
    \item принимает файлы расписания для импорта;
    \item извлекает из сообщения \texttt{chatId}, \texttt{telegramId},
    текст команды, callback data и fileId;
    \item передаёт распознанную команду в \texttt{MessageHandlerModule};
    \item отправляет пользователю ответы, уведомления и сообщения об ошибках.
\end{itemize}

\subsubsection{API (Healthcheck)}

\textbf{Роль:} служебный endpoint для проверки состояния системы.

\textbf{Функционал:}
\begin{itemize}
    \item обрабатывает запрос \texttt{GET /healthcheck};
    \item не требует административной авторизации;
    \item передаёт запрос в \texttt{HealthCheckModule};
    \item возвращает состояние сервера, состояние подключения к MongoDB и
    информацию об авторах проекта.
\end{itemize}

\subsubsection{API (Admin level)}

\textbf{Роль:} входной HTTP API для административных функций.

\textbf{Функционал:}
\begin{itemize}
    \item обрабатывает запрос \texttt{/auth} для получения токена администратора;
    \item обрабатывает запрос \texttt{/users} для получения списка пользователей;
    \item передаёт логин, пароль и Bearer-токен в \texttt{AdminAuthModule};
    \item передаёт авторизованные запросы в \texttt{ListUsersModule}.
\end{itemize}

\subsection{Модули системы}

\subsubsection{MessageHandlerModule}

\textbf{Роль:} центральный маршрутизатор пользовательских команд.

\textbf{Функционал:}
\begin{itemize}
    \item получает распознанную Telegram-команду из \texttt{API (Telegram message handlers)};
    \item определяет тип команды и выбирает нужный бизнес-модуль;
    \item направляет команды расписания в \texttt{TimetableModule};
    \item направляет команды импорта в \texttt{TimetableImportModule};
    \item направляет команды управления задачами в \texttt{TaskModule};
    \item направляет команды напоминаний в \texttt{ReminderModule};
    \item направляет запросы к LLM-советнику в \texttt{AdvisorModule};
    \item направляет команды синхронизации календаря в \texttt{CalendarSyncModule};
    \item направляет команды учёта учебного времени в \texttt{StudySessionModule};
    \item получает результат выполнения операции и формирует ответ пользователю.
\end{itemize}

\subsubsection{UserModule}

\textbf{Роль:} управление пользователями Telegram-бота.

\textbf{Функционал:}
\begin{itemize}
    \item обрабатывает команду \texttt{/start};
    \item идентифицирует пользователя по Telegram \texttt{chatId} или
    \texttt{telegramId};
    \item создаёт нового пользователя при первом обращении к боту;
    \item хранит настройки пользователя: имя, username, часовой пояс,
    настройки напоминаний;
    \item обращается к \texttt{UserRepository} для чтения и сохранения
    пользовательских данных.
\end{itemize}

\subsubsection{TimetableModule}

\textbf{Роль:} управление расписанием занятий и экзаменов.

\textbf{Функционал:}
\begin{itemize}
    \item выводит расписание на сегодня, завтра и текущую неделю;
    \item добавляет занятия и экзамены вручную;
    \item обновляет данные существующих занятий и экзаменов;
    \item удаляет события расписания;
    \item группирует события по дате и сортирует их по времени начала;
    \item обращается к \texttt{TimetableEventRepository} для работы с событиями
    расписания.
\end{itemize}

\subsubsection{TimetableImportModule}

\textbf{Роль:} импорт расписания из файлов и внешних источников.

\textbf{Функционал:}
\begin{itemize}
    \item принимает CSV- или iCal-файл, загруженный пользователем;
    \item выполняет проверку формата файла;
    \item разбирает строки расписания и преобразует их в события системы;
    \item валидирует дату, время, тип события, предмет и обязательные поля;
    \item может получать расписание из внешнего API учебного заведения через
    \texttt{API (Institution Timetable)};
    \item сохраняет корректные занятия и экзамены через
    \texttt{TimetableEventRepository};
    \item сохраняет результат импорта и ошибки через \texttt{ImportJobRepository}.
\end{itemize}

\subsubsection{TaskModule}

\textbf{Роль:} управление учебными задачами пользователя.

\textbf{Функционал:}
\begin{itemize}
    \item добавляет учебные задачи, дедлайны, упражнения, лабораторные и
    курсовые работы;
    \item связывает задачи с учебными предметами;
    \item выводит задачи на сегодня, завтра и неделю;
    \item выводит задачи по выбранному предмету;
    \item изменяет статус задачи на выполненную;
    \item обращается к \texttt{StudyTaskRepository} для сохранения и получения
    задач.
\end{itemize}

\subsubsection{ReminderModule}

\textbf{Роль:} управление напоминаниями.

\textbf{Функционал:}
\begin{itemize}
    \item хранит и изменяет настройки напоминаний пользователя;
    \item создаёт напоминания о занятиях, экзаменах и дедлайнах;
    \item определяет время отправки уведомлений;
    \item обращается к \texttt{ReminderRepository} для хранения настроек и
    задач напоминаний;
    \item публикует задачи на отправку уведомлений через
    \texttt{API (RabbitMQ)}.
\end{itemize}

\subsubsection{NotificationModule}

\textbf{Роль:} отправка уведомлений пользователю.

\textbf{Функционал:}
\begin{itemize}
    \item получает задачи уведомлений из RabbitMQ через \texttt{API (RabbitMQ)};
    \item формирует текст уведомления для пользователя;
    \item отправляет напоминания о занятиях, экзаменах и дедлайнах через
    \texttt{API (Telegram message handlers)};
    \item обрабатывает ошибки отправки сообщений.
\end{itemize}

\subsubsection{AdvisorModule}

\textbf{Роль:} LLM-советник по учебному планированию.

\textbf{Функционал:}
\begin{itemize}
    \item обрабатывает команды \texttt{/priority}, \texttt{/exam\_plan} и
    \texttt{/ask\_advisor};
    \item собирает контекст пользователя: расписание, экзамены, задачи,
    дедлайны и приоритеты;
    \item обращается к \texttt{TimetableEventRepository} и
    \texttt{StudyTaskRepository} для получения учебного контекста;
    \item формирует prompt для внешней LLM;
    \item отправляет запрос во внешний LLM API через \texttt{API (LLM)};
    \item возвращает пользователю учебные приоритеты, план подготовки к
    экзамену или ответ на свободный вопрос.
\end{itemize}

\subsubsection{CalendarSyncModule}

\textbf{Роль:} синхронизация расписания и задач с внешним календарём.

\textbf{Функционал:}
\begin{itemize}
    \item обрабатывает команды синхронизации календаря;
    \item передаёт занятия, экзамены и дедлайны во внешний календарь;
    \item работает с Google Calendar API или Yandex Calendar API через
    \texttt{API (Calendar)};
    \item сохраняет связь между локальными событиями и внешними событиями
    календаря через \texttt{CalendarSyncRepository};
    \item обрабатывает ошибки синхронизации.
\end{itemize}

\subsubsection{StudySessionModule}

\textbf{Роль:} учёт учебного времени пользователя.

\textbf{Функционал:}
\begin{itemize}
    \item обрабатывает команды \texttt{/study\_start}, \texttt{/study\_stop} и
    \texttt{/study\_stats};
    \item фиксирует начало и окончание учебной сессии;
    \item связывает учебную сессию с предметом;
    \item рассчитывает длительность учебной сессии;
    \item предоставляет статистику по предметам и периодам;
    \item обращается к \texttt{StudySessionRepository} для хранения истории
    занятий.
\end{itemize}

\subsubsection{AdminAuthModule}

\textbf{Роль:} аутентификация и авторизация администратора.

\textbf{Функционал:}
\begin{itemize}
    \item обрабатывает запросы к endpoint \texttt{/auth};
    \item проверяет логин и пароль администратора;
    \item обращается к \texttt{AdminRepository} для получения данных
    администратора;
    \item генерирует API-ключ или Bearer-токен при успешной аутентификации;
    \item проверяет токен при обращении к защищённым административным
    endpoint-ам.
\end{itemize}

\subsubsection{ListUsersModule}

\textbf{Роль:} получение списка пользователей для администратора.

\textbf{Функционал:}
\begin{itemize}
    \item обрабатывает административный запрос \texttt{/users};
    \item получает список пользователей через \texttt{UserRepository};
    \item возвращает администратору Telegram ID, username и другую служебную
    информацию о пользователях;
    \item использует \texttt{AdminAuthModule} для проверки прав доступа.
\end{itemize}

\subsubsection{HealthCheckModule}

\textbf{Роль:} проверка состояния системы.

\textbf{Функционал:}
\begin{itemize}
    \item обрабатывает запросы от \texttt{API (Healthcheck)};
    \item проверяет состояние серверного приложения;
    \item проверяет доступность MongoDB;
    \item при необходимости проверяет доступность RabbitMQ;
    \item возвращает статус системы и список авторов проекта.
\end{itemize}

\subsection{Репозитории}

\textbf{Роль:} слой доступа к данным MongoDB.

\textbf{Функционал:}
\begin{itemize}
    \item предоставляет бизнес-модулям операции чтения, записи, обновления и
    удаления данных;
    \item инкапсулирует работу с коллекциями MongoDB;
    \item передаёт запросы к базе данных через \texttt{API (MongoDB)}.
\end{itemize}

\subsubsection{UserRepository}

\textbf{Роль:} хранение данных пользователей.

\textbf{Функционал:}
\begin{itemize}
    \item сохраняет Telegram ID, chatId, username и настройки пользователя;
    \item предоставляет поиск пользователя по Telegram ID;
    \item используется модулями \texttt{UserModule} и \texttt{ListUsersModule}.
\end{itemize}

\subsubsection{AdminRepository}

\textbf{Роль:} хранение данных администраторов.

\textbf{Функционал:}
\begin{itemize}
    \item хранит логин администратора и хеш пароля;
    \item используется модулем \texttt{AdminAuthModule};
    \item не хранит пароль администратора в открытом виде.
\end{itemize}

\subsubsection{TimetableEventRepository}

\textbf{Роль:} хранение занятий и экзаменов.

\textbf{Функционал:}
\begin{itemize}
    \item хранит события расписания пользователя;
    \item различает типы событий: занятие и экзамен;
    \item предоставляет поиск событий на день, неделю и по предмету;
    \item используется модулями \texttt{TimetableModule},
    \texttt{TimetableImportModule} и \texttt{AdvisorModule}.
\end{itemize}

\subsubsection{StudyTaskRepository}

\textbf{Роль:} хранение учебных задач.

\textbf{Функционал:}
\begin{itemize}
    \item хранит задачи, дедлайны, приоритеты и статусы выполнения;
    \item предоставляет поиск задач по предмету и периоду;
    \item используется модулями \texttt{TaskModule} и \texttt{AdvisorModule}.
\end{itemize}

\subsubsection{ReminderRepository}

\textbf{Роль:} хранение настроек и задач напоминаний.

\textbf{Функционал:}
\begin{itemize}
    \item хранит настройки времени напоминаний;
    \item хранит информацию о запланированных уведомлениях;
    \item используется модулем \texttt{ReminderModule}.
\end{itemize}

\subsubsection{StudySessionRepository}

\textbf{Роль:} хранение учебных сессий.

\textbf{Функционал:}
\begin{itemize}
    \item хранит начало, окончание и длительность учебной сессии;
    \item связывает сессию с пользователем и предметом;
    \item используется модулем \texttt{StudySessionModule}.
\end{itemize}

\subsubsection{CalendarSyncRepository}

\textbf{Роль:} хранение информации о синхронизации с календарём.

\textbf{Функционал:}
\begin{itemize}
    \item хранит связь между локальным событием и внешним событием календаря;
    \item хранит статус последней синхронизации;
    \item используется модулем \texttt{CalendarSyncModule}.
\end{itemize}

\subsubsection{ImportJobRepository}

\textbf{Роль:} хранение истории импорта расписания.

\textbf{Функционал:}
\begin{itemize}
    \item сохраняет статус импорта расписания;
    \item хранит количество успешно добавленных и пропущенных записей;
    \item хранит ошибки импорта;
    \item используется модулем \texttt{TimetableImportModule}.
\end{itemize}

\subsection{External API adapters}

\textbf{Роль:} адаптеры для взаимодействия с внешними системами.

\textbf{Функционал:}
\begin{itemize}
    \item скрывают детали конкретных внешних API от бизнес-модулей;
    \item выполняют сериализацию и десериализацию запросов;
    \item обрабатывают ошибки внешних сервисов;
    \item возвращают бизнес-модулям данные в едином формате.
\end{itemize}

\subsubsection{API (MongoDB)}

\textbf{Роль:} адаптер доступа к MongoDB.

\textbf{Функционал:}
\begin{itemize}
    \item обеспечивает подключение к MongoDB;
    \item выполняет CRUD-операции с документами;
    \item используется всеми репозиториями системы.
\end{itemize}

\subsubsection{API (RabbitMQ)}

\textbf{Роль:} адаптер взаимодействия с RabbitMQ.

\textbf{Функционал:}
\begin{itemize}
    \item публикует сообщения о запланированных уведомлениях;
    \item читает сообщения из очередей;
    \item передаёт задачи уведомлений в \texttt{NotificationModule};
    \item обеспечивает асинхронную обработку напоминаний.
\end{itemize}

\subsubsection{API (LLM)}

\textbf{Роль:} адаптер взаимодействия с внешним LLM API.

\textbf{Функционал:}
\begin{itemize}
    \item принимает prompt от \texttt{AdvisorModule};
    \item отправляет запрос во внешний LLM API;
    \item получает ответ модели;
    \item возвращает учебные рекомендации, план подготовки или ответ на
    свободный вопрос пользователя.
\end{itemize}

\subsubsection{API (Calendar)}

\textbf{Роль:} адаптер взаимодействия с внешними календарями.

\textbf{Функционал:}
\begin{itemize}
    \item создаёт события во внешнем календаре;
    \item обновляет и удаляет события при изменении локального расписания;
    \item получает внешние идентификаторы календарных событий;
    \item используется модулем \texttt{CalendarSyncModule}.
\end{itemize}

\subsubsection{API (Institution Timetable)}

\textbf{Роль:} адаптер взаимодействия с API расписания учебного заведения.

\textbf{Функционал:}
\begin{itemize}
    \item отправляет запросы во внешний API расписания;
    \item получает список занятий и экзаменов;
    \item передаёт полученные данные в \texttt{TimetableImportModule};
    \item используется как один из источников импорта расписания.
\end{itemize}

\subsection{MongoDB}

\textbf{Роль:} основная база данных системы.

\textbf{Функционал:}
\begin{itemize}
    \item хранит пользователей;
    \item хранит занятия и экзамены;
    \item хранит учебные задачи и дедлайны;
    \item хранит настройки напоминаний;
    \item хранит учебные сессии;
    \item хранит данные синхронизации с внешними календарями;
    \item хранит историю импортов расписания;
    \item хранит данные администраторов.
\end{itemize}

\subsection{RabbitMQ}

\textbf{Роль:} брокер сообщений для асинхронной обработки задач.

\textbf{Функционал:}
\begin{itemize}
    \item принимает сообщения о необходимости отправки напоминаний;
    \item хранит задачи уведомлений до момента обработки;
    \item передаёт сообщения в \texttt{NotificationModule};
    \item повышает надёжность отправки уведомлений и отделяет логику
    планирования от логики отправки.
\end{itemize}

\subsection{LLM API}

\textbf{Роль:} внешний сервис генерации учебных рекомендаций.

\textbf{Функционал:}
\begin{itemize}
    \item принимает prompt с учебным контекстом пользователя;
    \item генерирует рекомендации по приоритетам;
    \item генерирует план подготовки к экзамену;
    \item отвечает на свободные вопросы пользователя о планировании учёбы.
\end{itemize}

\subsection{Google/Yandex Calendar API}

\textbf{Роль:} внешний календарный сервис.

\textbf{Функционал:}
\begin{itemize}
    \item принимает запросы на создание календарных событий;
    \item обновляет события при изменении расписания или дедлайнов;
    \item удаляет события при отключении синхронизации;
    \item возвращает внешние идентификаторы созданных событий.
\end{itemize}

\subsection{Institution Timetable API}

\textbf{Роль:} внешний источник расписания учебного заведения.

\textbf{Функционал:}
\begin{itemize}
    \item предоставляет расписание занятий и экзаменов;
    \item может использовать номер группы, направление или другой идентификатор
    учебной группы;
    \item возвращает данные, которые затем преобразуются в события локального
    расписания пользователя.
\end{itemize}

\subsection{Потоки данных на диаграмме}

Пользователь отправляет Telegram-команды, сообщения и файлы в
\texttt{API (Telegram message handlers)}. Далее запрос передаётся в
\texttt{MessageHandlerModule}, который маршрутизирует его в нужный модуль.

При работе с расписанием \texttt{MessageHandlerModule} вызывает
\texttt{TimetableModule}. Модуль получает или изменяет занятия и экзамены
через \texttt{TimetableEventRepository}, который выполняет операции с MongoDB
через \texttt{API (MongoDB)}.

При импорте расписания \texttt{TimetableImportModule} получает файл CSV/iCal
или данные из \texttt{Institution Timetable API}. После валидации данные
сохраняются через \texttt{TimetableEventRepository}. Результаты импорта и
ошибки сохраняются через \texttt{ImportJobRepository}.

При работе с задачами \texttt{TaskModule} сохраняет и получает задачи через
\texttt{StudyTaskRepository}. Задачи могут использоваться не только для вывода
списков, но и как часть контекста для \texttt{AdvisorModule}.

При настройке напоминаний \texttt{ReminderModule} сохраняет параметры через
\texttt{ReminderRepository} и публикует задачу уведомления через
\texttt{API (RabbitMQ)}. Затем \texttt{NotificationModule} получает задачу
уведомления и отправляет сообщение пользователю через Telegram.

При обращении к LLM-советнику \texttt{AdvisorModule} собирает расписание,
экзамены, задачи и дедлайны пользователя, формирует prompt и отправляет его
через \texttt{API (LLM)} во внешний \texttt{LLM API}. Ответ модели возвращается
пользователю через Telegram.

При синхронизации календаря \texttt{CalendarSyncModule} передаёт занятия,
экзамены и дедлайны через \texttt{API (Calendar)} во внешний
\texttt{Google/Yandex Calendar API}. Связь между локальными и внешними
событиями сохраняется через \texttt{CalendarSyncRepository}.

Администратор обращается к \texttt{API (Admin level)}. Для получения доступа
используется endpoint \texttt{/auth}, а для получения списка пользователей --
endpoint \texttt{/users}. Проверка данных администратора выполняется в
\texttt{AdminAuthModule}, а список пользователей получается через
\texttt{ListUsersModule} и \texttt{UserRepository}.

Endpoint \texttt{/healthcheck} обрабатывается через \texttt{API (Healthcheck)}
и \texttt{HealthCheckModule}. Он возвращает состояние сервера, состояние
подключения к MongoDB и служебную информацию о проекте.
\newpage
\begin{landscape}
\begin{figure}[H]
\thispagestyle{empty}
    \centering
    \includegraphics[scale=0.5]{arc.png}
    \caption{Архитектура}
    \vspace{30mm}17
    \label{arc}
\end{figure}
\end{landscape}
\end{document}